\begin{table}[htbp]
\centering
\footnotesize
\caption{Definizione alternativa di bene ambientale: lista OCSE contro lista APEC}
\label{tab:apec}
\begin{threeparttable}
\begin{tabular}{lcc}
\toprule
 & Lista OCSE & Lista APEC \\
 & (usata nel lavoro) & (54 codici) \\
\midrule
Profondit\`a EP (WB) $\times$ Verde & -0.00457 & 0.00501 \\
 & (0.00696) & (0.01266) \\
 & [0.512] & [0.693] \\
\addlinespace
Conteggio EP (TREND) $\times$ Verde & 0.00181 & 0.00319 \\
 & (0.00182) & (0.00208) \\
 & [0.320] & [0.126] \\
\bottomrule
\end{tabular}
\begin{tablenotes}[flushleft]\footnotesize
\item Variante baseline, pannello collassato.
\item \textbf{A cosa serve.} Non esiste una definizione unica di ``bene ambientale''. Il lavoro usa la lista OCSE; qui la si sostituisce con la lista APEC, molto pi\`u ristretta e concordata politicamente nel 2012. Se il risultato dipendesse da quali prodotti chiamiamo verdi, le due colonne divergerebbero.
\item Errore standard fra parentesi tonde, $p$-value fra parentesi quadre. $^{*}$ $p<0.10$, $^{**}$ $p<0.05$, $^{***}$ $p<0.01$.
\end{tablenotes}
\end{threeparttable}
\end{table}
